% HAND-WRITTEN fragment -- NOT generated by maestro2tex. Input from
% AnalogChapter.tex immediately after AmpAB.tex, which is generated.
%
% Salvaged from the retired note_calibration_seq.tex steps 4 and 5 and
% note_calibration_paths.tex C2, plus the fabric of
% ~/vestarv/docs/afe_rev2_debug_mux.html: per-channel slot 5 = ce_amp (A1
% output, inboard of the C1 isolation switch), slot 4 = vout, AFE_ADCSRCSEL = 4
% routes ATB into the SAR (§6.3). The "do not buffer it" argument is §4.1's
% chicken-and-egg: a test buffer built from the same 2.5 V devices carries an
% offset of the same class as A1 (sigma 5.20 mV) and A2 (sigma 8.91 mV).
% Numbers: tab_ampab_mc_offset, tab_tiaab_mc_offset. No cell names in prose.

\subsubsection{Measuring this block on chip}

Input offset is the term that decides yield, and a loopback is the only way to
see it. Open the reference-electrode isolation switch, close the unity-gain
loopback around A1, select per-channel slot 5 --- the amplifier output, inboard
of the counter-electrode isolation switch --- and read it either on
\texttt{atb\_d} with \texttt{AFE\_DBGDIREN}~=~1 or on-chip by routing the bus
into the converter with \texttt{AFE\_ADCSRCSEL}~=~4 and averaging. Repeat at
three reference codes across the working window: the intercept is the input
offset, \(\sigma = \SI{5.20}{\milli\volt}\)
(Table~\ref{tab:ampab-mc-offset}), and the slope is the finite-loop-gain
tracking error. Take the reading unbuffered --- a test buffer built from the
same devices carries an offset of the same class, so buffering substitutes one
unknown for another. For the transimpedance role, isolate the working electrode
and read slot 4 the same way; that reading is output-referred and carries
\(\sigma = \SI{8.91}{\milli\volt}\) (Table~\ref{tab:tiaab-mc-offset}), which
divided by the active gain is the channel's zero-current error. This test fabric
is design intent; it is not yet in the schematics.
